% 第四章 通信接口与协议设计

\section{通信接口与协议设计}

\subsection{设备间通信设计}
\par 本系统的设备间通信主要包含两条UART链路，分别用于“采集数据上报”和“视觉控制交互”。相较于复杂的二进制帧格式，本文在工程验证阶段采用了可读性更强的文本协议，并通过换行符作为一条消息的边界，结合超时清空缓冲区的方式实现容错。该设计便于串口调试与日志分析，并能在数据量较小的家庭环境监测场景中满足实时性要求。
\par STM32与ESP32-S3之间用于环境数据上报的UART链路采用按行上报策略。采集侧将各传感器值以键值对形式拼接并以逗号分隔，行末以换行符结束。网关侧按字符接收并累积到缓冲区，当检测到换行符时将当前行作为一帧数据进行解析；当缓冲区在一定时间内未收到新字符时，网关侧清空缓冲区以避免半包造成的解析错误。该策略在“低带宽、小消息、高可读性”的场景下具有实现简单、可维护性强的优势。
\par ESP32-S3与OpenMV之间用于控制与响应的UART链路采用“命令—响应”模式。网关侧向OpenMV发送命令字符串（如采集、识别、停止、状态查询、图像获取、名单管理等），OpenMV侧解析命令并执行相应任务后回传文本响应或JPEG图像数据。为兼顾文本与图像两类数据，网关侧在接收端实现了模式识别机制：当检测到JPEG文件头时切换到二进制接收模式，直到检测到JPEG结束标记或超时；随后将图像数据进行Base64编码，并以统一的文本载荷在MQTT链路中传输。对于多行文本响应，网关侧通过短超时窗口判断“响应结束”，将多行内容合并为一次响应，避免拆分导致的上层解析复杂度上升。
\par 错误处理方面，串口链路采用“超时丢弃与重试”为主的策略。上报链路中，采集数据周期性发送，单帧丢失不会影响系统整体趋势分析；控制链路中，网关侧通过超时机制判定OpenMV响应是否到达，并在必要时向上层返回超时信息。对于耗时较长的采集与识别任务，网关侧会延长等待窗口并允许进度消息持续回传，从而避免误判超时。

\par 配图建议如下。图4-1建议为“UART数据帧格式示意图”，画法为数据帧结构图：展示起始字符、键值对字段（T/H/P/VOC）、逗号分隔符、换行符结束，并在下方标注字段含义与示例数据。图4-2建议为“UART接收与解析流程示意图”，画法为流程图：字符接收→累积缓冲区→检测换行符→解析键值对→更新数据结构→超时清空，并在关键分支标注处理逻辑。
\begin{figure}[H]
\centering
\fbox{\parbox{0.92\linewidth}{\centering 图4-1占位：UART数据帧格式示意图。\\
要求展示帧结构（键值对+逗号分隔+换行符），并给出示例数据行（如T:25.3,H:60.5,P:1013.2,VOC:125,）。}}
\caption{UART数据帧格式示意图（建议配图）}
\label{fig:uart_frame}
\end{figure}
\begin{figure}[H]
\centering
\fbox{\parbox{0.92\linewidth}{\centering 图4-2占位：UART接收与解析流程示意图。\\
要求展示从字符接收到数据解析的完整流程，标注换行符检测、超时清空等关键分支。}}
\caption{UART接收与解析流程示意图（建议配图）}
\label{fig:uart_parsing}
\end{figure}

\subsection{网络通信设计}
\par 网络通信设计包含两条主要链路，分别用于数据上报与控制交互。数据上报链路采用WiFi + HTTP方式，将网关侧融合后的环境数据以JSON格式周期上传至服务器。控制交互链路采用MQTT发布/订阅模型，在服务器与网关之间建立命令下发与响应回传通道，适合处理异步事件与多消息交互场景\cite{mqtt_oasis}。
\par WiFi连接管理采用非阻塞策略。网关在启动后发起异步连接，并在主循环中轮询连接状态；当连接成功时记录网络状态并启动后续网络业务；当连接超时或运行过程中断开时，网关按固定间隔自动重连并更新状态。该设计避免了在网络不稳定时阻塞主循环，从而保障串口接收、显示刷新与控制响应的连续性。
\par HTTP数据上报采用REST风格接口。网关将温度、湿度、气压、VOC与气体传感器读数等字段组织为JSON对象，并附带本地时间戳以辅助端侧调试；服务器在接收到数据后添加服务器时间戳并写入数据存储。该方式实现简单，便于与Web前端共享同一套数据接口。
\par MQTT控制链路采用固定主题设计以降低复杂度。网关侧订阅命令主题，接收来自服务器的控制指令；网关侧向响应主题发布OpenMV的执行结果或图像数据。对于图像数据，由于Base64编码后数据量较大，网关侧在超过缓冲区限制时采用分块发布策略，并在主题上增加分块子主题以避免单条消息超过客户端缓冲区。服务器侧在接收响应时区分文本与图像载荷，并将识别结果等异步消息存储为可查询的历史记录。
\par 网络异常处理以“可恢复”为目标。MQTT侧启用自动重连机制并在断线期间保持命令队列的可控处理；当命令超时或响应丢失时，服务器向前端返回明确的错误信息，提示用户重试或检查连接。通过在网关与服务器两侧同时记录连接状态与关键日志，可以较快定位问题并提升系统可维护性。

\par 配图建议如下。图4-3建议为“HTTP数据上报流程时序图”，画法为时序图：ESP32采集数据→构建JSON→HTTP POST请求→服务器接收→添加时间戳→写入存储→返回响应，并在关键步骤标注时间戳添加点。图4-4建议为“MQTT主题与消息流向图”，画法为发布/订阅模型图：展示命令主题（smarthome/command）与响应主题（smarthome/response），标注服务器发布命令、网关订阅命令、网关发布响应、服务器订阅响应的流向。图4-5建议为“图像Base64编码与分块传输流程示意图”，画法为流程图：JPEG二进制→Base64编码→检查大小→超过阈值则分块→MQTT分块发布→服务器拼装→前端渲染，并在分块步骤标注块头格式（CHUNK:index/total:data）。
\begin{figure}[H]
\centering
\fbox{\parbox{0.92\linewidth}{\centering 图4-3占位：HTTP数据上报流程时序图。\\
要求展示从ESP32采集到服务器写入的完整时序，标注JSON构建、POST请求、时间戳添加等关键步骤。}}
\caption{HTTP数据上报流程时序图（建议配图）}
\label{fig:http_sequence}
\end{figure}
\begin{figure}[H]
\centering
\fbox{\parbox{0.92\linewidth}{\centering 图4-4占位：MQTT主题与消息流向图。\\
要求展示命令主题与响应主题，标注发布/订阅关系与消息流向（服务器→网关→OpenMV→网关→服务器）。}}
\caption{MQTT主题与消息流向图（建议配图）}
\label{fig:mqtt_topics}
\end{figure}
\begin{figure}[H]
\centering
\fbox{\parbox{0.92\linewidth}{\centering 图4-5占位：图像Base64编码与分块传输流程示意图。\\
要求展示JPEG→Base64→大小检查→分块发布→服务器拼装→前端渲染的完整流程，标注分块头格式。}}
\caption{图像Base64编码与分块传输流程示意图（建议配图）}
\label{fig:image_transmission}
\end{figure}
